\medskip

\emph{Principaux domaines abordés}: Manipulation des vecteurs, des droites et des plans de l'espace. Orthogonalité et distances dans l'espace. Représentations paramétriques et équations cartésiennes.

\bigskip

Le solide ABCDEFGH est un cube. On se place dans le repère orthonormé $\left(\text{A}~;~\vect{\imath},~\vect{\jmath},~\vect{k}\right)$ de l'espace dans lequel les coordonnées des points B, D et E sont : 

\[\text{B}(3~;~0 ~;~0) , \text{D} (0~;~3~;~0)\: \text{et E}(0~;~0~;~3).\]

\begin{center}
\psset{unit=1cm}
\begin{pspicture}(6.5,5.8)
\pspolygon(0.5,1)(2.9,0.5)(2.9,3.7)(0.5,4.2)
\psline(2.9,0.5)(5.6,1)(5.6,4.2)(2.9,3.7)
\psline(5.6,4.2)(3.2,4.7)(0.5,4.2)
\psline[linestyle=dashed](0.5,1)(3.2,1.5)(5.6,1)
\psline[linestyle=dashed](3.2,1.5)(3.2,4.7)
\psline[linewidth=1.3pt]{->}(0.5,1)(1.3,0.84)
\psline[linewidth=1.3pt]{->}(0.5,1)(1.4,1.16)
\psline[linewidth=1.3pt]{->}(0.5,1)(0.5,2.06)
\uput[dl](0.5,1){A} \uput[d](2.9,0.5){B} \uput[r](5.6,1){C} \uput[d](3.2,1.5){D}
\uput[l](0.5,4.2){E} \uput[u](2.9,3.7){F} \uput[r](5.6,4.2){G} \uput[u](3.2,4.7){H}
\uput[d](0.7,0.9){$\vect{\imath}$} \uput[u](0.95,1.08){$\vect{\jmath}$} \uput[l](0.5,1.53){$\vect{k}$}
\end{pspicture}
\end{center}

On considère les points P(0~;~0~;~1), Q(0~;~2~;~3) et R(1~;~0~;~3).

\medskip

\begin{enumerate}
\item Placer les points P, Q et R sur la figure.
\begin{center}


\psset{unit=1.5cm}
\begin{pspicture}(6.5,5.8)
\pspolygon(0.5,1)(2.9,0.5)(2.9,3.7)(0.5,4.2)
\psline(2.9,0.5)(5.6,1)(5.6,4.2)(2.9,3.7)
\psline(5.6,4.2)(3.2,4.7)(0.5,4.2)
\psline[linestyle=dashed](0.5,1)(3.2,1.5)(5.6,1)
\psline[linestyle=dashed](3.2,1.5)(3.2,4.7)
\psline[linewidth=1.3pt]{->}(0.5,1)(1.3,0.84)
\psline[linewidth=1.3pt]{->}(0.5,1)(1.4,1.16)
\psline[linewidth=1.3pt]{->}(0.5,1)(0.5,2.06)
\uput[dl](0.5,1){A} \uput[d](2.9,0.5){B} \uput[r](5.6,1){C} \uput[d](3.2,1.5){D}
\uput[l](0.5,4.2){E} \uput[u](2.9,3.7){F} \uput[r](5.6,4.2){G} \uput[u](3.2,4.7){H}
\uput[d](0.7,0.9){$\vect{\imath}$} \uput[u](0.95,1.08){$\vect{\jmath}$} \uput[l](0.5,1.53){$\vect{k}$}


\end{pspicture}
\end{center}

\item Montrer que le triangle PQR est isocèle en R.
\item Justifier que les points P{}, Q et R définissent un plan.
\item On s'intéresse à présent à la distance entre le point E et le plan (PQR).
	\begin{enumerate}
		\item Montrer que le vecteur $\vect{u}(2~;~1~;~- 1)$ est normal au plan (PQR).
		\item En déduire une équation cartésienne du plan (PQR).
		\item Déterminer une représentation paramétrique de la droite $(d)$ passant par le point E et orthogonale au plan (PQR).
		\item Montrer que le point L$\left(\dfrac23~;~\dfrac13~;~\dfrac83\right)$ est le projeté orthogonal du point E sur le plan (PQR).
		\item Déterminer la distance entre le point E et le plan (PQR).
	\end{enumerate}
\item En choisissant le triangle EQR comme base, montrer que le volume du tétraèdre EPQR est $\dfrac23$.

On rappelle que le volume $V$ d'un tétraèdre est donné par la formule : 

\[V = \dfrac13 \times \text{aire d'une base} \times \text{hauteur correspondante}.\]

\item Trouver, à l'aide des deux questions précédentes, l'aire du triangle PQR.
\end{enumerate}

\bigskip
